% !TEX root = ../../../main.tex
\section{Uniform convergence}
\label{sec:analysis-i:uniform-convergence}

% BEGIN TUFTE PLACEHOLDER: supplied sample, lines 1077--1098.
\begingroup\sloppy
% Replace this entire specimen block with reviewed course notes.
\newthought{The font and style} of the marginal material may also be modified using the following commands:

\begin{docspec}
  \doccmd{setsidenotefont}\{\docopt{font commands}\}\\
  \doccmd{setcaptionfont}\{\docopt{font commands}\}\\
  \doccmd{setmarginnotefont}\{\docopt{font commands}\}\\
  \doccmd{setcitationfont}\{\docopt{font commands}\}
\end{docspec}

The \doccmddef{setsidenotefont} sets the font and style for sidenotes, the
\doccmddef{setcaptionfont} for captions, the \doccmddef{setmarginnotefont} for
margin notes, and the \doccmddef{setcitationfont} for citations.  The
\docopt{font commands} can contain font size changes (e.g.,
\doccmdnoindex{footnotesize}, \doccmdnoindex{Huge}, etc.), font style changes (e.g.,
\doccmdnoindex{sffamily}, \doccmdnoindex{ttfamily}, \doccmdnoindex{itshape}, etc.), color changes (e.g.,
\doccmdnoindex{color}\texttt{\{blue\}}), and many other adjustments.

If, for example, you wanted the captions to be set in italic sans serif, you could use:
\begin{docspec}
  \doccmd{setcaptionfont}\{\doccmdnoindex{itshape}\doccmdnoindex{sffamily}\}
\end{docspec}

\lipsum[3]
\par\endgroup
% END TUFTE PLACEHOLDER
